\documentclass[10pt]{article}
\providecommand{\FIGDIR}{fig}
\newcommand{\DIGESTDATE}{12 August 2026}
% ===================== PM Research Watch - shared preamble =====================
\usepackage[T1]{fontenc}
\usepackage[utf8]{inputenc}
\usepackage{amsmath,amssymb}
\usepackage{tgpagella}
\usepackage{tgheros}
\renewcommand{\sfdefault}{qhv}
\renewcommand{\ttdefault}{lmtt}
\usepackage{microtype}
\usepackage[english]{babel}

\usepackage[a4paper,top=17mm,bottom=17mm,left=17mm,right=17mm,
            headsep=6mm,footskip=10mm]{geometry}
\usepackage{graphicx}
\usepackage{xcolor}
\usepackage{colortbl}
\usepackage{tikz}
\usetikzlibrary{calc}
\usepackage[most]{tcolorbox}
\usepackage{titlesec}
\usepackage{fancyhdr}
\usepackage{booktabs}
\usepackage{array}
\usepackage{longtable}
\usepackage{enumitem}
\usepackage{multicol}
\usepackage{ragged2e}
\usepackage[hidelinks]{hyperref}

% ------------------------------------------------------------------ palette
\definecolor{Ink}    {HTML}{14262E}
\definecolor{Deep}   {HTML}{0F4C5C}
\definecolor{Teal}   {HTML}{1F7A8C}
\definecolor{Sky}    {HTML}{4FA3C4}
\definecolor{Sage}   {HTML}{6FA287}
\definecolor{Amber}  {HTML}{E3A008}
\definecolor{Clay}   {HTML}{B5643C}
\definecolor{Coral}  {HTML}{D1495B}
\definecolor{Violet} {HTML}{6A5B9E}
\definecolor{Slate}  {HTML}{4A6572}
\definecolor{Paper}  {HTML}{FBF9F5}
\definecolor{Mist}   {HTML}{EEEAE2}
\definecolor{Rule}   {HTML}{D9D3C9}

\pagecolor{Paper}
\color{Ink}

\hypersetup{colorlinks=true, linkcolor=Deep, urlcolor=Teal, citecolor=Teal,
            pdftitle={PM Research Watch}, pdfauthor={Automated daily digest}}

% ------------------------------------------------------------------ headings
\titleformat{\section}
  {\sffamily\bfseries\large\color{Deep}}{}{0pt}
  {}[{\vspace{1.2mm}\color{Amber}\titlerule[1.1pt]}]
\titlespacing*{\section}{0pt}{5.5mm}{2.8mm}

\titleformat{\subsection}
  {\sffamily\bfseries\normalsize\color{Teal}}{}{0pt}{}
\titlespacing*{\subsection}{0pt}{3.6mm}{1.4mm}

\setlist[itemize]{leftmargin=4.2mm, itemsep=1.1mm, topsep=1.3mm, parsep=0pt,
                  label=\textcolor{Teal}{\footnotesize$\blacktriangleright$}}

% ------------------------------------------------------------------ header/footer
\pagestyle{fancy}
\fancyhf{}
\renewcommand{\headrulewidth}{0pt}
\renewcommand{\footrulewidth}{0pt}
\fancyhead[L]{\sffamily\scriptsize\color{Slate}PM RESEARCH WATCH\;\textcolor{Rule}{\textbar}\;\DIGESTDATE}
\fancyhead[R]{\sffamily\scriptsize\color{Slate}Particulate matter monitoring \& health effects}
\fancyfoot[C]{\sffamily\scriptsize\color{Slate}\thepage}

% ------------------------------------------------------------------ components
% metric tile
\newcommand{\metric}[3]{%
  \begin{tcolorbox}[enhanced, colback=white, colframe=Rule, boxrule=0.5pt,
    arc=1.6mm, left=2.2mm, right=2.2mm, top=1.8mm, bottom=1.8mm,
    borderline west={1.7mm}{0pt}{#3}, nobeforeafter, valign=center,
    height=17mm, width=\linewidth]
    \RaggedRight
    {\sffamily\bfseries\fontsize{17}{17}\selectfont\color{#3}#1}\\[0.5mm]
    {\sffamily\fontsize{6.5}{7.4}\selectfont\color{Slate}#2}
  \end{tcolorbox}}

% pull-quote / key insight box
\newtcolorbox{keybox}[1][]{enhanced, colback=Mist, colframe=Mist,
  arc=1.6mm, left=3.2mm, right=3.2mm, top=2.4mm, bottom=2.4mm,
  borderline west={1.7mm}{0pt}{Amber}, #1}

% subtopic band
\newcommand{\band}[2]{%
  \vspace{3.2mm}\noindent
  \begin{tcolorbox}[enhanced, colback=#2, colframe=#2, arc=1.2mm,
    left=2.6mm, right=2.6mm, top=1.1mm, bottom=1.1mm, boxrule=0pt,
    width=\linewidth, nobeforeafter]
    {\sffamily\bfseries\small\color{white}#1}
  \end{tcolorbox}\vspace{1.6mm}}

% one paper entry: {short cite}{title}{journal}{doi}{body}
\newcommand{\paperentry}[5]{%
  \par\vspace{1.9mm}\noindent
  {\sffamily\bfseries\footnotesize\color{Ink}#1}\;%
  {\sffamily\fontsize{7}{8}\selectfont\color{Slate}#3}\\[0.4mm]
  {\footnotesize\itshape\color{Deep}#2}\\[0.6mm]
  {\footnotesize #4}\\[0.3mm]
  {\fontsize{6.6}{7.6}\selectfont\color{Slate}\texttt{doi:}\,\href{https://doi.org/#5}{\texttt{\detokenize{#5}}}}
  \par}

% caption style
\newcommand{\figcap}[1]{\par\vspace{1.0mm}
  {\fontsize{7.2}{8.6}\selectfont\color{Slate}\textsf{\textbf{Fig.}}\; #1\par}\vspace{1.5mm}}

\providecommand{\permil}{\textperthousand}

\setlength{\parindent}{0pt}
\setlength{\parskip}{0pt}
\renewcommand{\arraystretch}{1.18}


\begin{document}

% ==================================================================== MASTHEAD
\thispagestyle{empty}
\begin{tikzpicture}[remember picture, overlay]
  \fill[Deep] (current page.north west) rectangle
      ($(current page.north east)+(0,-67mm)$);
  \fill[Amber] ($(current page.north west)+(0,-67mm)$) rectangle
      ($(current page.north east)+(0,-68.4mm)$);
  \foreach \x/\y/\r in {22/12/0.9, 41/26/1.5, 63/9/0.7, 88/31/1.1, 112/17/1.9,
                        138/29/0.8, 159/13/1.3, 178/54/1.0, 196/21/1.6,
                        30/44/0.6, 74/57/1.2, 125/48/0.9, 168/6/0.8, 191/59/0.7,
                        52/18/1.0, 100/60/0.6, 148/20/1.1, 66/49/1.4, 205/45/0.8}
    \fill[Sky, opacity=0.30] ($(current page.north west)+(\x mm,-\y mm)$) circle (\r mm);
\end{tikzpicture}

\vspace*{4mm}
{\sffamily\fontsize{7.6}{9}\selectfont\color{Sky}%
 AUTOMATED DAILY LITERATURE DIGEST\;\textbullet\;PhD RESEARCH WATCH\par}
\vspace{2.2mm}
{\sffamily\bfseries\fontsize{27}{29}\selectfont\color{white}%
 Particulate Matter\par}
{\sffamily\fontsize{27}{29}\selectfont\color{Sky}%
 Monitoring \& Health Effects\par}
\vspace{5.0mm}
{\sffamily\fontsize{8.4}{10}\selectfont\color{white}%
 Issue 2026-08-12\;\;\textcolor{Amber}{\textbar}\;\;Entry window 12 August 2026\;\;%
 \textcolor{Amber}{\textbar}\;\;28 records carried, 14 rejected\par}
\vspace{18mm}

% -------------------------------------------------------------- metric strip
\noindent
\begin{minipage}[t]{0.190\linewidth}\metric{28}{RECORDS CARRIED\\(14 REJECTED)}{Deep}\end{minipage}\hfill
\begin{minipage}[t]{0.190\linewidth}\metric{8}{SUBTOPIC CLUSTERS\\REPRESENTED}{Teal}\end{minipage}\hfill
\begin{minipage}[t]{0.190\linewidth}\metric{5}{EFFECT ESTIMATES\\WITH INTERVALS}{Sky}\end{minipage}\hfill
\begin{minipage}[t]{0.190\linewidth}\metric{11}{INSTRUMENTATION /\\EXPOSURE RECORDS}{Amber}\end{minipage}\hfill
\begin{minipage}[t]{0.190\linewidth}\metric{6}{METADATA-ONLY\\(NO ABSTRACT)}{Clay}\end{minipage}

\vspace{6mm}

% -------------------------------------------------------------- signal of day
\section*{\sffamily\bfseries\color{Deep}Signal of the day}
\vspace{-1.5mm}{\color{Amber}\rule{\linewidth}{1.1pt}}\vspace{2.5mm}

\begin{keybox}
{\sffamily\bfseries\small\color{Deep}The Parkinson's--PM$_{2.5}$ association is carried by
exposure at least four years before the event --- which makes the exposure-history
requirement, not the concentration, the binding constraint on the design.}\\[1.4mm]
\footnotesize
Delaney et al.\ (\textit{Environ Epidemiol}) fitted distributed-lag discrete-time survival
models to \textbf{10{,}366{,}083 Medicare fee-for-service beneficiaries} aged 65+ across the
contiguous US, 2000--2016, with 10-year exposure histories for PM$_{2.5}$, NO$_2$ and
summer ozone. Ten years of continuous PM$_{2.5}$ at the 90th versus the 0.5th percentile
(11.8 vs.\ 3.0\,$\mu$g/m$^3$) gave \textbf{OR 1.634 (95\% CI 1.489--1.792)} for first
hospitalisation with PD; NO$_2$ at 31.7 vs.\ 3.7\,ppb gave \textbf{1.474 (1.379--1.575)}.
Ozone was equivocal. The load-bearing detail is the lag structure: associations were
carried by exposure \emph{at least four years} before hospitalisation, with the
near-event years contributing little.

Two consequences follow for instrumentation. First, the exposure contrast that matters
here is well below current US design values --- 11.8\,$\mu$g/m$^3$ is an unremarkable
annual mean --- so the measurement problem is long-run bias and drift, not peak capture.
A network whose calibration wanders by 1--2\,$\mu$g/m$^3$ over a multi-year deployment
degrades exactly the contrast this model exploits. Second, hospitalisation with PD is a
proxy for onset that is displaced by an unknown and probably differential interval; the
four-year lag should be read as a lower bound on latency, not an estimate of it.
\end{keybox}

\vspace{3mm}

% -------------------------------------------------------------- exec brief
\section*{\sffamily\bfseries\color{Deep}What changed today}
\vspace{-1.5mm}{\color{Amber}\rule{\linewidth}{1.1pt}}\vspace{2.5mm}

\footnotesize
\begin{itemize}
\item \textbf{Emission-factor choice is invisible in total mass and decisive in source
attribution.} Park et al.\ ran EMEP/EEA and KIMM non-exhaust factors through the same
3-D CFD--chemistry model over a dense Seoul district. Total non-exhaust PM$_{10}$ was
similar under both; the partitioning was not --- brake wear dominant under EMEP/EEA,
road wear dominant under KIMM. Non-exhaust reached \textbf{up to 70\% of total
PM$_{10}$} in near-road cells. A monitor reporting mass alone cannot distinguish the two
worlds, and they imply opposite mitigation priorities.

\item \textbf{Two independent mediation results converge on oxidative/inflammatory
metabolism.} Chow et al.\ (\textit{Nat Commun}) profiled $>$1{,}100 pre-diagnostic plasma
metabolites in 1{,}357 Cancer Prevention Study participants and found eight metabolites
associated with both exposure and subsequent lung cancer; four
($\gamma$-glutamylglutamine, phenylacetylglutamate, N-(2-furoyl)glycine,
4-vinylguaiacol glucuronide) mediated the PM and ozone associations at $q<0.2$. On the
same day, a correspondence in \textit{JACC Adv} pushes back on the statistical fragility
of exactly this class of metabolomic mediation analysis --- worth reading as a pair.

\item \textbf{Renal burden is now quantified with a pooled exposure--response function.}
Li et al.\ pooled 38 CKD and 16 eGFR studies: \textbf{RR 1.32 (1.21--1.43) per
10\,$\mu$g/m$^3$}, predominantly log-linear, with a saturating eGFR response at high
concentration. Attributed 2020 burden is 485{,}000 deaths and 243 million prevalent
cases; PM$_{2.5}$ ranks second among level-2 CKD mortality risk factors, up from first in
1990 only because metabolic risks grew faster.

\item \textbf{Instrumentation: a non-flammable CPC working fluid reaches parity.}
Kirchhoff et al.\ ran a DMSO-based condensation particle counter against a butanol
counterpart for six months. Mean cutoff \textbf{5.8$\pm$0.9\,nm} (DMSO) vs.\
\textbf{5.6$\pm$0.5\,nm} (butanol), with substantially lower fluid consumption. The
practical unlock is airborne and unattended remote deployment, where butanol's
flammability is the blocking constraint.

\item \textbf{Open-source exposure tooling.} Chillrud et al.\ released \texttt{pcpr},
an R implementation of principal component pursuit with three environmental-health
extensions (non-negativity, missingness, below-LOD penalty). Applied to 26 PM$_{2.5}$
constituents in Queens 2015--2021 it recovered four patterns and 211 outlying events
including July Fourth firework spikes. Directly usable on speciated network output.

\item \textbf{The abstract famine continues --- fifth consecutive issue.} Six records
carried today have no abstract at the entry date: four Elsevier
\textit{Atmos.\ Environ.} articles, one Elsevier correspondence, one Lancet editorial.
The four \textit{Atmos.\ Environ.} papers are squarely in scope (mobile monitoring,
S\~ao Paulo PM$_{10}$ speciation, Australian biomass-burning fingerprints, absorbing
aerosol index) and carry no extractable content. They are queued for the Saturday
weekly sweep.
\end{itemize}

\clearpage

% ==================================================================== ANALYTICS
\section*{\sffamily\bfseries\color{Deep}Corpus analytics}
\vspace{-1.5mm}{\color{Amber}\rule{\linewidth}{1.1pt}}\vspace{3mm}

\noindent\includegraphics[width=\linewidth]{\FIGDIR/f1_subtopics.png}
\figcap{Subtopic assignment of the 28 carried records, single-label by dominant
contribution. Exposure assessment and sensing together account for 11 records (39\%),
the highest instrumentation share since 3 August; the balance is unusually
review-heavy on the health side. Neither \textit{Reproductive \& developmental} nor
\textit{Occupational \& indoor} received a record today.}

\vspace{2mm}
\noindent\includegraphics[width=\linewidth]{\FIGDIR/f2_design_endpoint.png}
\figcap{Left --- study architecture. The \textit{Metadata only} slice (6) is reported as
its own category rather than pooled with characterised designs; it is a property of the
harvest, not of the science. Right --- health endpoint. Twelve records carry no human
health endpoint at all, which is what an instrumentation-heavy day looks like.}

\vspace{2mm}
\noindent\includegraphics[width=\linewidth]{\FIGDIR/f3_metric_geography.png}
\figcap{Left --- particle metric. PM$_{2.5}$ dominates, but note three records whose
exposure metric is not a mass concentration at all (aerosol optical depth, absorbing
aerosol index, ultrafine particle number) --- these are the records that a mass-only
monitoring programme cannot produce. Right --- geographic provenance. No record today
originates from sub-Saharan Africa, and the single South Asian record is a bench device
study.}

\vspace{2mm}
\noindent\includegraphics[width=\linewidth]{\FIGDIR/f6_lifecourse.png}
\figcap{Life-course windows addressed. Counts are non-exclusive and describe this issue
only. Both the in-utero and pubertal windows are empty today, an unusual gap given the
last fortnight's developmental-origins run; the mass sits in working-age and 65+
cohorts.}

\vspace{2mm}
\noindent\includegraphics[width=\linewidth]{\FIGDIR/f4_heatmap.png}
\figcap{Subtopic $\times$ study-architecture cross-tabulation. The structural feature to
note is the same one the 27 July issue flagged and which has not moved: the
\textit{Sensing} and \textit{Exposure assessment} rows contain no health-endpoint
linkage, and the health rows contain no instrumentation. The two literatures are
running in parallel within this harvest.}

\vspace{2mm}
\noindent\includegraphics[width=\linewidth]{\FIGDIR/f5_forest.png}
\figcap{The five interval estimates carried today, on a common logarithmic axis.
\textbf{Exposure contrasts are heterogeneous and the estimates are not mutually
comparable}: Delaney's two ORs are 90th-vs-0.5th-percentile contrasts over a 10-year
history, Li's RR is per 10\,$\mu$g/m$^3$, Jain's is a pooled highest-vs-lowest RR across
studies with $I^2=78.7\%$, and Schlaepfer's is a weighted cumulative first-year index.
The plot shows direction and precision only. Nothing here is pooled. Kakhaia et al.'s
counterfactual RR of 0.90 is omitted: the abstract reports a confidence interval on the
risk difference and none on the ratio, and the projected-policy interval as printed
($-1.05$; CI $-1.12$ to $-1.07$) does not contain its own point estimate.}

\clearpage

% ==================================================================== DIGEST
\section*{\sffamily\bfseries\color{Deep}Paper-level digest}
\vspace{-1.5mm}{\color{Amber}\rule{\linewidth}{1.1pt}}\vspace{1mm}

{\fontsize{7.4}{9}\selectfont\color{Slate}Ordered by cluster. Each entry gives the
methodological core and the finding that matters, not an abstract paraphrase.\par}

\band{Neuro, cognition \& mental health\hfill 1 record}{Deep}

\paperentry{Delaney et al.\ 2026}
{Cumulative annual air pollution exposure and first hospitalisation with Parkinson's disease}
{Environ Epidemiol}
{Discrete-time survival with distributed lag models over 10{,}366{,}083 Medicare FFS
beneficiaries 65+, contiguous US 2000--2016; 10-year exposure histories for PM$_{2.5}$,
NO$_2$ and summer O$_3$; non-linear exposure--response. \textbf{PM$_{2.5}$ OR 1.634
(1.489--1.792)} at 11.8 vs.\ 3.0\,$\mu$g/m$^3$; \textbf{NO$_2$ OR 1.474 (1.379--1.575)}
at 31.7 vs.\ 3.7\,ppb; O$_3$ limited. Associations carried by exposure $\geq$4 years
pre-event. The endpoint is hospitalisation, not incidence --- diagnosis-to-hospitalisation
lag is unmeasured and plausibly differential by exposure via healthcare access, which is
the main threat to the latency claim. PM$_{2.5}$ and NO$_2$ are strongly correlated in
Medicare-scale data and the paper does not resolve which is doing the work.}
{10.1097/EE9.0000000000000519}

\band{Cardiovascular \& metabolic\hfill 1 record}{Teal}

\paperentry{Kakhaia et al.\ 2026}
{Exposome approach to causal effects of time-varying air pollution mixtures on CVD incidence}
{Environ Int}
{UK Biobank, annual PM$_{2.5}$ and NO$_2$ from a Europe-wide LUR model, 10-year
cumulative CVD incidence under modified treatment policies. Confounders selected by a
causal structure discovery algorithm with prior knowledge; sequential doubly robust
non-parametric estimation with ensemble learning and censoring adjustment. A 25\%
exposure reduction over 2005--2019 gave $-0.62$ per 100 (95\% CI $-0.70$ to $-0.55$),
RR\,$\approx$\,0.90; the projected net-zero scenario reached $-1.05$ per 100 and
RR\,$\approx$\,0.85. \textbf{The stated sensitivity is the finding}: dropping
assessment-centre region from the adjustment set more than doubled the estimate, so the
headline number is contingent on a spatial-confounding choice the authors cannot
adjudicate. Method-forward, estimate-fragile.}
{10.1016/j.envint.2026.110441}

\band{Other clinical endpoints\hfill 5 records}{Slate}

\paperentry{Chow et al.\ 2026}
{Blood metabolomic signatures linking air pollution to lung cancer}
{Nat Commun}
{Prospective nested case-control across CPS-II and CPS-3, 1{,}357 participants,
$>$1{,}100 metabolites in pre-diagnostic plasma, residential concentrations of six
pollutants at blood draw. Eight metabolites associated with both exposure and subsequent
lung cancer; four ($\gamma$-glutamylglutamine, phenylacetylglutamate,
N-(2-furoyl)glycine, 4-vinylguaiacol glucuronide) significantly mediated PM and ozone
associations at adjusted $q<0.2$. Single-timepoint metabolomics cannot separate exposure
response from early disease, and $q<0.2$ is a permissive threshold over 1{,}100 features
--- treat the four as candidates. Smoking adjustment in a lung cancer metabolome study is
the obvious residual-confounding exposure.}
{10.1038/s41467-026-75116-3}

\paperentry{Li et al.\ 2026}
{Global and national CKD burden attributable to fine particles: updated exposure--response meta-analysis}
{Innovation (Camb)}
{Three-stage mixed-effects meta-analysis over 38 CKD and 16 eGFR studies, fitting a
non-linear ERF and propagating it to 1990--2020 burden above a 10\,$\mu$g/m$^3$
counterfactual. \textbf{RR 1.32 (1.21--1.43) per 10\,$\mu$g/m$^3$}, predominantly
log-linear; \textbf{eGFR $-2.115$ (0.452--3.777)\,mL/min/1.73\,m$^2$} with saturation at
high exposure. 2020 burden: 485{,}000 deaths, 5.77 million incident and 243 million
prevalent cases. The counterfactual floor is a choice, not a threshold, and the burden
figures scale with it; the primary studies are also concentrated in high-exposure Asian
settings, so the low-concentration arm of the ERF is thin.}
{10.1016/j.xinn.2026.101429}

\paperentry{Jain et al.\ 2026}
{Ambient air pollution and glaucoma: systematic review and meta-analysis}
{J Curr Glaucoma Pract}
{PRISMA review, PubMed and Embase to 30 December 2025; 18 observational studies, 8
pooled by generic inverse variance. \textbf{Pooled RR 1.12 (1.06--1.18)} for long-term
PM$_{2.5}$, $I^2=78.7\%$. NO$_2$, NO$_x$, PM$_{10}$ and CO directionally positive but
synthesised narratively. RNFL and GCIPL thinning and modest IOP elevation also reported.
Heterogeneity at 79\% with eight studies makes the pooled point estimate a weak summary;
the authors say causality cannot be established and they are right. Two databases only
is a real recall limitation for this literature.}
{10.5005/jp-journals-10078-1522}

\paperentry{Schlaepfer et al.\ 2026}
{Environmental hazards and the risk of idiopathic urethral stricture}
{Transl Androl Urol}
{Utah Population Database, 2{,}048 cases born 1996--2021 matched 1:1 on birth year and
location to 1{,}939 controls; historical air pollution and toxics exposure assessed
prenatally, in the first year of life, and in the five years before diagnosis.
\textbf{Only the first-year window was associated: RR 1.21 (1.07--1.36).} Univariate
signals for ethylbenzene, tetrachloroethylene and molybdenum trioxide. Three exposure
windows tested with one positive is a multiplicity problem the paper does not address,
and the univariate toxicant hits are unadjusted. Included as a hypothesis, not a result.}
{10.21037/tau-2026-0338}

\paperentry{Liu et al.\ 2026}
{Reply: statistical fragility and methodological constraints in metabolomic mediation}
{JACC Adv}
{Correspondence from the Emory group (Liu, Chang, Jones, Quyyumi, Liang) responding to a
critique of metabolomic mediation between air pollution and cardiovascular risk.
\textbf{Abstract not indexed} --- carried as metadata only because it sits directly
against Chow et al.\ (same senior author, same method class) in today's harvest and the
methodological dispute is the reason to read both. Full text required.}
{10.1016/j.jacadv.2026.103149}

\clearpage

\band{Respiratory \& allergic\hfill 4 records}{Sage}

\paperentry{Feng \& Chao 2026}
{Lower respiratory infection burden attributable to PM$_{2.5}$ in China and the G20, 1990--2023}
{Int J Biometeorol}
{GBD 2023 deaths and DALYs from LRI attributable to ambient and household PM$_{2.5}$,
Joinpoint regression on age-standardised DALY rates with sex subgroups and loess against
SDI. Long decline reversed upward in 2021--2023 (APC \textbf{China 12.76--14.23\%}, G20
\textbf{14.94--16.04\%} for ambient); household-attributable rates kept falling in China
but rose in the G20 from 2021 (APC 8.49--11.20\%). Burden concentrated in under-5s and
$\geq$70s. The post-COVID reversal is the interesting claim and it rests on three
data-years of GBD modelled estimates, which are themselves smoothed --- an inflection
this recent should be treated as provisional.}
{10.1007/s00484-026-03283-5}

\paperentry{Yang et al.\ 2026}
{COPD and LRI burden attributable to household air pollution and secondhand smoke, 1990--2021}
{J Thorac Dis}
{GBD 2021 across 204 countries with BAPC and ARIMA projection. Age-standardised COPD
DALYs and mortality fell 37.00\% and 37.12\%; LRI fell 66.35\% and 53.62\%. Burden stays
concentrated in LMICs, in older women for COPD and in under-15s and $\geq$60s for LRI;
HAP-attributable COPD declined faster than SHS-attributable. Projected female COPD
burden from both sources exceeds male by 2030. Standard GBD secondary analysis --- the
exposure is a modelled solid-fuel-use proxy, not measured PM, and the projections inherit
every assumption in the parent estimates.}
{10.21037/jtd-2026-0284}

\paperentry{Wu et al.\ 2026}
{Air pollution and COPD-related health outcomes: bibliometric and visualisation analysis}
{J Thorac Dis}
{11{,}167 Web of Science articles and reviews (1964 to 31 May 2025) through CiteSpace,
VOSviewer and \texttt{bibliometrix}: country/institution collaboration, journal
co-citation, keyword clustering, burst detection. USA, China and UK lead output; Harvard
most influential institution; China and climate change identified as 2021-onward
frontiers. Bibliometrics measure publication behaviour, not evidence --- useful only for
orienting a literature search or a grant narrative.}
{10.21037/jtd-2026-0784}

\paperentry{The Lancet Respiratory Medicine 2026}
{The far-reaching effects of wildfire smoke}
{Lancet Respir Med}
{Unsigned editorial. \textbf{Abstract not indexed}; carried as metadata only. Logged
because wildfire-source attribution has been the recurring through-line of this watch
since 27 July and an editorial in this journal signals where the clinical field is
placing it. Full text required for any citable content.}
{10.1016/S2213-2600(26)00256-0}

\band{Mechanistic toxicology\hfill 2 records}{Violet}

\paperentry{Choi et al.\ 2026}
{Degradable porous PLGA/PCL membrane enabling a lung alveoli-on-a-chip for particulate injury}
{Biofabrication}
{Porogen-assisted phase-separation membranes with tunable degradation ($\sim$40\% surface
porosity gain, $\sim$3\,$\mu$m thickness loss over culture) carrying primary human
alveolar epithelium and microvascular endothelium at an air--liquid interface under
breathing-mimetic cyclic strain. Apical diesel PM produced dose-dependent cytotoxicity,
barrier permeability increase, elevated ROS and DNA damage \emph{in both layers} ---
trans-barrier propagation, which an insert-only model cannot show. Roflumilast-N-oxide
(PDE4 inhibitor) selectively attenuated oxidative and inflammatory readouts without
restoring barrier integrity. Deposited dose is delivered apically in suspension and is
not equivalent to inhaled dosimetry; no dose--response anchoring to ambient
concentrations is attempted. Relevant here as a bench receptor for source-resolved
particle samples.}
{10.1088/1758-5090/ae9402}

\paperentry{Yuan et al.\ 2026}
{Hallmarks of lung cancer driven by inhalable particulate matter}
{Innovation (Camb)}
{Critical review proposing a 12-feature ``hallmarks of PM-associated lung cancer''
framework spanning cell-intrinsic traits (genetic alteration, epigenetic reprogramming,
metabolic remodelling, stemness) and microenvironment remodelling (inflammation, immune
suppression, angiogenesis, pre-metastatic niche). Argues for integrating PM exposure into
screening and risk stratification. A framework paper: the hallmark taxonomy is an
organising device, not evidence, and the underlying studies are heterogeneous in particle
source and dose. Useful as a citation scaffold for a mechanism section.}
{10.1016/j.xinn.2026.101347}

\band{Burden, policy \& mitigation\hfill 4 records}{Coral}

\paperentry{Tian et al.\ 2026}
{Bioinspired honeycomb-structured nanofibrous membranes for high-efficiency PM capture}
{Nanomaterials (Basel)}
{Template-assisted electrospinning onto a mesh receiver establishes a directional
electric field, producing concentrated ordered stacking plus sparse random orientation of
fine-diameter fibres. \textbf{$>$98.51\% filtration efficiency at 34\,Pa pressure drop},
85\% transparency, 130.8\,mm/s air permeability. Efficiency is quoted for a single
challenge condition; no loading curve, no electrostatic-decay ageing, no humidity
sensitivity --- the three things that determine whether an electret-assisted filter holds
its number in service. Bench result, not a service life.}
{10.3390/nano16150941}

\paperentry{Xie et al.\ 2026}
{Thirty-year change and projection of Asian tracheal, bronchus and lung cancer burden}
{J Thorac Dis}
{GBD 2021 across Asian subregions with Joinpoint, age-period-cohort, decomposition,
frontier, inequality and BAPC prediction. Incidence roughly doubled and deaths rose
$\sim$2.7-fold 1990--2021; ASIR +22\%, ASPR +49\%, ASMR +8\%. Population growth was the
largest contributor, then ageing. \textbf{Particulate matter pollution ranks second
behind smoking among risk factors at $\sim$10\% and rising.} Six analytic frameworks on
one modelled dataset multiply presentation, not information.}
{10.21037/jtd-2026-0740}

\paperentry{Liu et al.\ 2026b}
{Lung cancer incidence and mortality trends with age-period-cohort analysis, Zibo 2014--2024}
{Front Oncol}
{Registry data for a heavily industrial Shandong city. Age-standardised incidence rose
(AAPC \textbf{+4.34\%}, 95\% CI 2.14--6.59) while mortality fell (AAPC \textbf{$-$2.13\%},
$-$3.53 to $-$0.72); male rates exceeded female throughout. Period effect increasing,
cohort effect decreasing across successive birth cohorts. Rising incidence with falling
mortality in a screening-active decade is the classic overdiagnosis signature and the
paper does not test it; air pollution appears as narrative context with no exposure
measurement at all.}
{10.3389/fonc.2026.1873592}

\paperentry{Suryanarayanareddy et al.\ 2026}
{Low-cost bio-based particle-capturing device for vehicular emissions: pilot study}
{Environ Sci Pollut Res Int}
{Cylindrical coconut-husk and sugarcane-pulp enclosure tapering to a conical outlet and
a monitoring chamber, tested indoors under simulated junction conditions at three
turbulence intensities. Single-pass removal 15--20\% (PM$_{2.5}$) and 15--45\%
(PM$_{10}$) for halted vehicles; 20--65\% and 15--65\% for moving vehicles. Efficiency
ordering flips with turbulence regime, which the paper reports honestly. Indoor
simulation of an outdoor near-road environment, no replication count given, and no
material-degradation testing on a biodegradable enclosure meant to sit at a roadside.}
{10.1007/s11356-026-38118-8}

\clearpage

\band{Exposure assessment \& modelling\hfill 6 records}{Clay}

\paperentry{Park et al.\ 2026}
{City-scale non-exhaust PM$_{10}$ from emissions to exposure with a coupled chemistry--CFD framework}
{Environ Pollut}
{Two non-exhaust emission-factor frameworks (EMEP/EEA Guidebook and KIMM) implemented in
the Fine-scale Air Quality Simulation Unit, a 3-D coupled CFD--chemistry model, over a
densely built Seoul district, combined with the 2022 Seoul Time-Activity Profile.
\textbf{Total non-exhaust PM$_{10}$ was similar under both frameworks; source
partitioning was not} --- brake wear dominant under EMEP/EEA, road wear under KIMM.
Non-exhaust reached \textbf{up to 70\% of total PM$_{10}$} in near-road and dense cells;
outdoor activity-weighted contributions of 9--10\% in commercial and traffic zones
exceeded concentration-based estimates because peak traffic coincides with peak outdoor
presence. No measurement validation of the partitioning is presented, which is the
obvious next step and the reason a receptor-model or tracer campaign would be
complementary.}
{10.1016/j.envpol.2026.128941}

\paperentry{Chillrud et al.\ 2026}
{\texttt{pcpr}: principal component pursuit for exposure pattern recognition and outlier detection}
{Environ Epidemiol}
{R package implementing PCP with three environmental-health extensions: a non-negativity
constraint on the low-rank matrix, missingness handling, and a specialised penalty for
below-LOD observations. Demonstrated on 26 PM$_{2.5}$ constituents measured every few
days in Queens, NYC, 2015--2021: four patterns recovered (secondary/tailpipe, traffic,
crustal dust, salt) plus \textbf{211 outlying events} including July Fourth firework
spikes. PCP's sparse component is an outlier detector, not a source; calling the 211
events ``exposure events'' presumes an interpretation the decomposition does not supply.
Still the most immediately usable artefact in today's corpus for anyone holding speciated
network data.}
{10.1097/EE9.0000000000000515}

\paperentry{Zhang et al.\ 2026}
{National-scale heavy metal contamination in urban road dust across 82 Chinese cities}
{Environ Geochem Health}
{13{,}269 samples pooled from 169 studies into an urban-agglomeration framework;
geo-accumulation index for anthropogenic enrichment, potential ecological risk index,
multivariate source deconvolution. Cd and Hg are the priority pollutants, jointly
\textbf{$>$85\% of total ecological risk} in every agglomeration despite low absolute
concentrations relative to Mn and Zn; a Cd/Hg $I_{\text{geo}}$ ratio is proposed as a
source fingerprint. Pooling 169 studies with heterogeneous sampling depth, particle-size
cut and digestion protocol is the central weakness --- the fingerprint may partly index
methodology. Road dust is the resuspension reservoir behind non-exhaust PM, so this pairs
with Park et al.\ above.}
{10.1007/s10653-026-03422-2}

\paperentry{Flor\^encio et al.\ 2026}
{Comprehensive analysis of PM$_{10}$ in the Metropolitan Area of S\~ao Paulo, Part I:
chemical characterisation and sources}
{Atmos Environ}
{\textbf{Abstract not deposited at Crossref on the entry date.} Title indicates
chemical speciation and source apportionment for a megacity with almost no
component-resolved PM literature in this watch to date. Queued for the Saturday sweep.}
{10.1016/j.atmosenv.2026.122286}

\paperentry{Stevens et al.\ 2026}
{Identifying unique biomass burning marker fingerprints for Australian hardwoods,
softwoods and grasses}
{Atmos Environ}
{\textbf{Abstract not deposited at Crossref on the entry date.} Fuel-specific tracer
fingerprints are the missing input for wildfire source attribution outside the Northern
Hemisphere; pairs with P\'erez-Pastor et al.\ below on the metrology side. Queued for
the Saturday sweep.}
{10.1016/j.atmosenv.2026.122303}

\paperentry{Ibrahim 2026}
{Air pollution status and temporal patterns in Bukit Rambai and Muar, 2013--2017}
{Research Square (preprint)}
{Five years of Malaysian Department of Environment station data (API, O$_3$, CO, SO$_2$,
PM$_{10}$, NO$_2$) at two southern Peninsular stations, analysed in \texttt{openair} with
calendar plots. Peak API 477 (Bukit Rambai) and \textbf{762 (Muar)} in June 2013 during a
transboundary haze episode. Descriptive, two stations, no exposure or health linkage, and
not peer reviewed. Carried only because episodic Southeast Asian haze is otherwise
unrepresented in this watch and the 762 value is a useful anchor for dynamic-range
requirements on low-cost optical sensors, which saturate well below it.}
{10.21203/rs.3.rs-10666765/v1}

\clearpage

\band{Sensing, forecasting \& instrumentation\hfill 5 records}{Amber}

\paperentry{Kirchhoff et al.\ 2026}
{Evaluation of DMSO as working fluid in condensation particle counters}
{Atmos Meas Tech}
{Laboratory and field comparison of a DMSO-operated CPC against a butanol counterpart
across pressure, temperature and aerosol type, with a modified automatic refill system
run for six months. Growth depends strongly on saturator temperature and
saturator--condenser $\Delta T$, confirmed by downstream optical particle counting and
saturation/droplet-size simulation. \textbf{Mean cutoff 5.8$\pm$0.9\,nm (DMSO) vs.\
5.6$\pm$0.5\,nm (butanol)}, with substantially lower fluid consumption; DMSO--H$_2$O
mixtures extend the operating range. The wider cutoff spread on the DMSO unit is the
honest cost, and matters if you are differencing two size channels. Non-flammability is
what makes airborne and unattended remote deployment tractable.}
{10.5194/amt-19-5243-2026}

\paperentry{P\'erez-Pastor et al.\ 2026}
{Robust quantification of anhydrosugars and reference material evaluation}
{Environ Sci Pollut Res Int}
{Accelerated solvent extraction (dichloromethane:methanol), BSTFA derivatisation, GC/MS
for levoglucosan, galactosan and mannosan. Linearity $R>0.9987$, repeatability $<$6\%,
LOD $<$0.02\,mg\,L$^{-1}$; relative error 1--7.4\% against SRM 2787 (fine PM).
\textbf{SRM 1649b (urban dust) gave extraction-condition-dependent levoglucosan values},
which the authors attribute to matrix ageing --- a direct warning against using an old
urban-dust SRM to validate a biomass-burning tracer method. Chromatographic calibration
dominates the uncertainty budget. Expanded uncertainties meet forthcoming EU air quality
monitoring requirements.}
{10.1007/s11356-026-38136-6}

\paperentry{Zhang et al.\ 2026b}
{Impact of aerosol-type assumption on Landsat 8 atmospheric correction over land}
{Atmos Meas Tech}
{600 Landsat 8 scenes over 100 AERONET sites, comparing the Urban Clean aerosol type
hardwired into the LaSRC surface-reflectance algorithm against MOD04- and MOD09-based
dynamic types. \textbf{MOD04 gave the best AOD accuracy: $R^2=0.7236$, RMSE 0.0437, bias
0.0052.} Relevant here not for surface reflectance but because the same fixed-aerosol-type
assumption propagates into satellite AOD products used as predictors in
satellite--surface PM$_{2.5}$ fusion models --- an $R^2$ of 0.72 against AERONET is the
ceiling any such predictor inherits. AERONET sites are systematically non-representative
of high-PM urban environments, so this is likely optimistic.}
{10.5194/amt-19-5281-2026}

\paperentry{Cao et al.\ 2026}
{Characterising on-road NO$_2$ and particulate matter across urban road environments
using city-scale mobile monitoring}
{Atmos Environ}
{\textbf{Abstract not deposited at Crossref on the entry date.} The most directly
relevant record of the day for mobile and distributed sensing and the one whose loss to
the abstract famine costs most. Queued for the Saturday sweep.}
{10.1016/j.atmosenv.2026.122302}

\paperentry{Zhang et al.\ 2026c}
{Characterisation of air pollution transport using absorbing aerosol index from multiple
payloads}
{Atmos Environ}
{\textbf{Abstract not deposited at Crossref on the entry date.} Multi-payload AAI
synergy is a transport-attribution method rather than a surface-concentration one;
logged for completeness. Queued for the Saturday sweep.}
{10.1016/j.atmosenv.2026.122298}

\clearpage

% ==================================================================== TABLE
\section*{\sffamily\bfseries\color{Deep}Corpus register}
\vspace{-1.5mm}{\color{Amber}\rule{\linewidth}{1.1pt}}\vspace{2.5mm}

{\fontsize{6.7}{8.4}\selectfont
\setlength{\tabcolsep}{2.4pt}
\begin{longtable}{@{}>{\RaggedRight}p{25mm}>{\RaggedRight}p{41mm}>{\RaggedRight}p{27mm}%
>{\RaggedRight}p{23mm}>{\RaggedRight}p{18mm}>{\RaggedRight}p{15mm}@{}}
\toprule
\textsf{\textbf{First author}} & \textsf{\textbf{Focus}} & \textsf{\textbf{Journal}} &
\textsf{\textbf{Design}} & \textsf{\textbf{Metric}} & \textsf{\textbf{Region}}\\
\midrule
\endfirsthead
\toprule
\textsf{\textbf{First author}} & \textsf{\textbf{Focus}} & \textsf{\textbf{Journal}} &
\textsf{\textbf{Design}} & \textsf{\textbf{Metric}} & \textsf{\textbf{Region}}\\
\midrule
\endhead
\multicolumn{6}{@{}l@{}}{\cellcolor{Mist}\textsf{\textbf{\textcolor{Deep}{Neuro / mental health}}}}\\[0.6mm]
\href{https://doi.org/10.1001/jamanetworkopen.2026.25015}{Dales} & Wildfire-sourced PM$_{2.5}$ and ED visits for migraine and other primary headache & \textit{JAMA Netw Open} & Case-crossover & PM$_{2.5}$ (wildfire) & Canada\\
\href{https://doi.org/10.1093/ageing/afag215}{Glover} & Umbrella meta-review of urban environmental exposures and cognitive health & \textit{Age Ageing} & Umbrella review / meta-analysis & PM$_{2.5}$ & Global\\
\href{https://doi.org/10.1016/j.nsa.2026.107019}{Goossens} & In utero air pollution exposure and pre-/postnatal brain development & \textit{Neurosci Appl} & Narrative review & PM$_{2.5}$ / PM$_{10}$ / PM$_{\mathrm{coarse}}$ & Global\\
\href{https://doi.org/10.1016/j.isci.2026.116617}{Ma} & PM$_{2.5}$-induced neurotoxicity and depression: exposure risks and mechanisms & \textit{iScience} & Narrative review & PM$_{2.5}$ + components & Global\\
\href{https://doi.org/10.3760/cma.j.cn112150-20250804-00752}{Pan} & Maternal PM exposure and offspring neurodevelopment: effects and mechanisms & \textit{Zhonghua Yu Fang Yi Xue Za Zhi} & Narrative review & PM (bulk) & China\\
\addlinespace[0.9mm]
\multicolumn{6}{@{}l@{}}{\cellcolor{Mist}\textsf{\textbf{\textcolor{Teal}{Cardiovascular \& metabolic}}}}\\[0.6mm]
\href{https://doi.org/10.3760/cma.j.cn112338-20251120-00843}{Huang} & PM$_{2.5}$ and dyslipidaemia in rural Chinese children; dietary effect modification & \textit{Zhonghua Liu Xing Bing Xue Za Zhi} & Cross-sectional (multistage) & PM$_{2.5}$ & China\\
\href{https://doi.org/10.32394/pe/218966}{Kazmierska} & Air pollution and cardiovascular disease: review of current data & \textit{Przegl Epidemiol} & Narrative review & PM$_{2.5}$ / PM$_{10}$ & Global\\
\href{https://doi.org/10.1007/s00125-026-06814-2}{Xi} & Air pollution, built environment and temperature vs. prediabetes phenotypes & \textit{Diabetologia} & Prospective cohort & PM$_{2.5}$ & Germany\\
\href{https://doi.org/10.1161/JAHA.125.046917}{Yang} & Short-term PM$_{2.5}$ and hospitalisation/ED use in atrial fibrillation & \textit{J Am Heart Assoc} & Case-crossover & PM$_{2.5}$ (wildfire-influenced) & USA\\
\addlinespace[0.9mm]
\multicolumn{6}{@{}l@{}}{\cellcolor{Mist}\textsf{\textbf{\textcolor{Sky}{Reproductive \& developmental}}}}\\[0.6mm]
\href{https://doi.org/10.1016/j.jes.2025.10.043}{Chen} & Prenatal air pollution, fetal growth trajectory phenotypes and child neurodevelopment & \textit{J Environ Sci (China)} & Prospective cohort & PM$_{2.5}$ / PM$_{10}$ & China\\
\href{https://doi.org/10.1016/j.ecoenv.2026.120566}{Huang} & Joint prenatal bisphenol + air-pollutant mixture exposure and preschool IQ & \textit{Ecotoxicol Environ Saf} & Prospective cohort (mixtures) & PM$_{2.5}$ / PM$_{10}$ & China\\
\href{https://doi.org/10.1007/s40572-026-00552-8}{Peterson} & Ambient air pollution, PAHs and pubertal development: critical review & \textit{Curr Environ Health Rep} & Critical review & PM$_{2.5}$ components & Global\\
\href{https://doi.org/10.1016/j.envres.2026.125313}{Xie} & Long-term PM$_{2.5}$ and timing of pubertal development in boys and girls & \textit{Environ Res} & Prospective cohort & PM$_{2.5}$ & China\\
\addlinespace[0.9mm]
\multicolumn{6}{@{}l@{}}{\cellcolor{Mist}\textsf{\textbf{\textcolor{Sage}{Respiratory \& allergic}}}}\\[0.6mm]
\href{https://doi.org/10.1007/s00484-026-03289-z}{Bayar Muluk} & Climate change and otolaryngology: airway disease impacts and emission reduction & \textit{Int J Biometeorol} & Narrative review & PM$_{2.5}$ & Global\\
\href{https://doi.org/10.4103/mgr.MEDGASRES-D-25-00164}{Liao} & Short-term air pollution, respiratory symptoms and lung function in asthma & \textit{Med Gas Res} & Panel study & PM$_{2.5}$ / PM$_{10}$ & China\\
\href{https://doi.org/10.3390/epidemiologia7040097}{Sanchez de Prada} & Ambient PM and asthma-coded admissions in children 0-3 y, Spain & \textit{Epidemiologia} & Ecological panel & PM$_{2.5}$ / PM$_{10}$ & Spain\\
\href{https://doi.org/10.1016/j.jhazmat.2026.143076}{Yang} & Allergenic vs. toxic risk of airborne seed hairs compared with ambient PM$_{2.5}$ & \textit{J Hazard Mater} & Field measurement & PM$_{2.5}$ + bioaerosol & China\\
\addlinespace[0.9mm]
\multicolumn{6}{@{}l@{}}{\cellcolor{Mist}\textsf{\textbf{\textcolor{Violet}{Mechanistic toxicology}}}}\\[0.6mm]
\href{https://doi.org/10.1016/j.jes.2025.12.074}{Ding} & PM$_{2.5}$ promotes vascular calcification via miR-27b-5p/HES1 signalling & \textit{J Environ Sci (China)} & In vitro + organoid + murine & PM$_{2.5}$ & China\\
\href{https://doi.org/10.1007/s00360-026-01698-5}{Maraschi} & Osmoregulatory disruption and metal bioaccumulation from settleable atmospheric PM & \textit{J Comp Physiol B} & Controlled exposure (ecotox) & Settleable APM & Brazil\\
\href{https://doi.org/10.1016/j.jes.2025.11.018}{Wu} & Macrophage NLRP3 inflammasome in real-ambient PM-induced hepatic glucose dysregulation & \textit{J Environ Sci (China)} & Real-ambient murine + co-culture & Real-ambient PM & China\\
\addlinespace[0.9mm]
\multicolumn{6}{@{}l@{}}{\cellcolor{Mist}\textsf{\textbf{\textcolor{Clay}{Occupational \& indoor}}}}\\[0.6mm]
\href{https://doi.org/10.3389/fmedt.2026.1796828}{Chen} & Nested suction instrument to cut intraoperative surgical-smoke PM$_{2.5}$ & \textit{Front Med Technol} & Bench / device experiment & PM$_{2.5}$ (surgical smoke) & China\\
\href{https://doi.org/10.1093/annweh/wxag055}{Cole} & Determinants of respirable crystalline silica exposure in Australian tunnelling & \textit{Ann Work Expo Health} & Observational exposure study & Respirable dust / RCS & Australia\\
\href{https://doi.org/10.1016/j.jhazmat.2026.143081}{Ding} & Stage-resolved hair Zn-isotope dynamics as a firefighting combustion-exposure tracer & \textit{J Hazard Mater} & Repeated-measures biomarker & Combustion-derived PM & China\\
\href{https://doi.org/10.3390/jox16040119}{Ramsden} & Xenobiotic hazards in aircraft cabin air & \textit{J Xenobiot} & Review + risk economics & Ultrafine particles & Global\\
\addlinespace[0.9mm]
\multicolumn{6}{@{}l@{}}{\cellcolor{Mist}\textsf{\textbf{\textcolor{Amber}{Sensing, forecasting \& instrumentation}}}}\\[0.6mm]
\href{https://doi.org/10.1016/j.jenvman.2026.130515}{Shi} & Adaptive spatiotemporal graph transformer for multi-site multi-step PM$_{2.5}$ forecasting & \textit{J Environ Manage} & Deep-learning model & PM$_{2.5}$ & China / India\\
\href{https://doi.org/10.3390/bios16070380}{Tsuchiya} & Seasonal and indoor-outdoor characteristics of fluorescent aerosol particles in offices & \textit{Biosensors} & Real-time sensor campaign & Bioaerosol / coarse & Japan\\
\addlinespace[0.9mm]
\multicolumn{6}{@{}l@{}}{\cellcolor{Mist}\textsf{\textbf{\textcolor{Amber}{Exposure assessment \& modelling}}}}\\[0.6mm]
\href{https://doi.org/10.1029/2025GH001621}{Solorzano-Araque} & Multiscale exposure modelling: LOTOS-EUROS + mobility surveys + particle toxicity & \textit{GeoHealth} & CTM + mobility integration & PM$_{2.5}$ / PM$_{10}$ & Colombia\\
\addlinespace[0.9mm]
\multicolumn{6}{@{}l@{}}{\cellcolor{Mist}\textsf{\textbf{\textcolor{Coral}{Burden, policy \& mitigation}}}}\\[0.6mm]
\href{https://doi.org/10.1002/hsr2.72886}{Karamzad} & Ischaemic heart disease burden in MENA 1990-2021; PM attributable fraction 28.2\% & \textit{Health Sci Rep} & GBD secondary analysis & Ambient PM & MENA\\
\href{https://doi.org/10.1371/journal.pone.0347989}{Rees} & HEPA filtration to reduce respiratory infection in care homes & \textit{PLoS One} & RCT process evaluation & Filtered PM & UK\\
\href{https://doi.org/10.1016/j.jes.2026.01.009}{Shan} & Policy-driven emission reductions and post-pandemic rebound, high-resolution inventories & \textit{J Environ Sci (China)} & Emission inventory & PM emissions & China\\
\addlinespace[0.9mm]
\multicolumn{6}{@{}l@{}}{\cellcolor{Mist}\textsf{\textbf{\textcolor{Slate}{Other clinical endpoints}}}}\\[0.6mm]
\href{https://doi.org/10.3390/dj14070440}{Lupita} & PM$_{2.5}$, preventive-care capacity and caries/edentulism burden across the EU & \textit{Dent J} & Ecological panel & PM$_{2.5}$ & EU-27\\
\href{https://doi.org/10.1016/j.jdin.2026.06.008}{Park} & Long-term PM exposure and risk of androgenetic alopecia & \textit{JAAD Int} & Nationwide cohort & PM$_{2.5}$ / PM$_{10}$ & South Korea\\
\href{https://doi.org/10.3171/2026.3.PEDS2615}{Sun} & Environmental exposure and area deprivation vs. paediatric hydrocephalus outcomes & \textit{J Neurosurg Pediatr} & Registry cohort + geospatial & PM$_{2.5}$ (satellite-derived) & USA\\
\addlinespace[0.9mm]
\bottomrule
\end{longtable}}


\vspace{4mm}
\begin{keybox}
{\sffamily\bfseries\footnotesize\color{Deep}Provenance, screening and known gaps}\\[1.2mm]
{\fontsize{7.2}{8.8}\selectfont
\textbf{Window.} \texttt{2026-08-12 -> 2026-08-12}, contiguous with the 11 August issue;
no gap. PubMed is windowed on \texttt{[Date - Entry]}; the Crossref-by-ISSN leg is
windowed on Crossref \texttt{created}. These are different clocks and are not treated as
interchangeable.\\[0.8mm]
\textbf{Per-leg yield.} PubMed connector, two queries (health axis, instrumentation
axis): 19 distinct, of which 1 was returned by both --- \textbf{18 carried}. Local
harvester: \texttt{pubmed\_health} 13 (5 additional records, 3 carried),
\texttt{pubmed\_sensing} \textbf{0}, Europe PMC 2 (\textbf{2 carried}, one of them a
Research Square preprint), Crossref by ISSN 16 (\textbf{7 carried}). arXiv returned
145\,kB of relevance-ranked results with nothing on or after 5 August in scope.
\textbf{OpenAlex HTTP 429, 10th consecutive} --- the leg is dead and is not being retried
beyond the single scripted call. Consensus returned 20 with the top 3 visible on the free
tier; all three resolved to papers with no verifiable entry date in the window and were
rejected rather than back-dated. \textbf{Sixth consecutive issue with zero Consensus
carry.} Scholar Gateway's corpus ends May 2026 and cannot serve an August entry window;
it was not used to fill the metadata-only records.
ClinicalTrials.gov: no PM or air-pollution registration updated in the window;
\texttt{trials.json} unchanged.\\[0.8mm]
\textbf{Rejections (14).} Off-topic: dissolved black carbon photodegradation on hematite;
meat sustainability labels; six ACP cloud/climate papers (ice supersaturation, CCN
replenishment, Hunga ozone, aerosol--cloud in ECMWF, secondary ice production, sulphate
AMOC response); N$_2$O clumped isotopes; IASI water vapour; stratospheric aerosol lidar.
Out-of-axis: sulfur-VOC + OH kinetics (SOA precursor chemistry, no PM mass or health
link); residential greenness and preeclampsia (exposure is NDVI, no PM estimate
reported); TB recurrence in Asia (PM one determinant among many, and no DOI registered to
verify). All logged with reasons to \texttt{state/rejected.jsonl}.\\[0.8mm]
\textbf{Metadata-only (6).} Four Elsevier \textit{Atmos.\ Environ.} articles, one
Elsevier \textit{JACC Adv} correspondence, one \textit{Lancet Respir Med} editorial. These
carry a title, journal and DOI and nothing else; every claim about them in this issue is
about the title, and is marked as such. This is the fifth consecutive issue with a
metadata-only block and the backlog now runs to 08-08.\\[0.8mm]
\textbf{Labels.} Subtopic, design, particle-metric and geography labels are drawn only
from the vocabularies already mapped in \texttt{plots.py} and \texttt{mktable.py}; no new
leaf label was written this issue. Assignment is single-label by dominant contribution
and is a judgement, not a controlled vocabulary.\\[0.8mm]
Content derived from PubMed, Europe PMC and Crossref metadata; abstracts remain the
copyright of their respective publishers. All DOI links resolve to the publisher of
record.\par}
\end{keybox}

\end{document}
